\chapter{Carpometacarpal (CMC) Joint Arthritis}

\section*{Definition and Overview}
Carpometacarpal (CMC) joint arthritis is osteoarthritis of the joint at the base of the thumb, where the first metacarpal bone meets the trapezium bone of the wrist. It is also known as basal thumb arthritis or thumb base arthritis. This saddle-shaped joint allows the thumb to oppose and pinch, making it essential for gripping and fine manipulation, but it is also one of the most heavily loaded joints in the hand. Arthritis here produces pain, stiffness, and weakness that interfere with pinching, gripping, and twisting activities such as opening jars, turning keys, and writing. CMC arthritis is among the most common sites of osteoarthritis in the upper limb, particularly in post-menopausal women, and it is a leading cause of hand disability in older adults.

\section*{Epidemiology}
Osteoarthritis of the thumb base is very common and increases steadily with age. It is considerably more frequent in women, who are affected several times more often than men, with the peak incidence after the menopause. Radiographic evidence of CMC arthritis is present in a large proportion of older people, although many of these individuals have no symptoms at all. Recognised risk factors include a family history of osteoarthritis, previous thumb or wrist injury, occupations and hobbies that involve repetitive pinching or gripping, and generalised joint hypermobility. Symptoms most often begin in middle age and progress gradually over many years.

\section*{Aetiology and Pathophysiology}
CMC arthritis is a degenerative (wear-and-tear) condition in which the articular cartilage of the thumb base joint is progressively lost, exposing underlying bone and leading to pain, stiffness, and deformity.
\begin{itemize}
    \item \textbf{Primary osteoarthritis:} with age, cartilage thins, roughens, and fragments, so that bone rubs against bone during movement.
    \item \textbf{Previous injury:} a fracture of the thumb metacarpal or scaphoid, or a severe sprain, can damage the joint surface and accelerate arthritis.
    \item \textbf{Repetitive loading:} decades of forceful pinching, gripping, and twisting load the joint beyond what its cartilage can tolerate.
    \item \textbf{Hormonal factors:} the strong female preponderance after the menopause suggests a hormonal influence on cartilage health.
    \item \textbf{Ligamentous laxity:} laxity of the anterior oblique ligament allows subluxation of the joint, concentrating wear on the palmar cartilage.
    \item \textbf{Genetic factors:} a family history of hand or generalised osteoarthritis increases susceptibility.
\end{itemize}

\section*{Clinical Features}
Symptoms typically develop insidiously and are centred on the base of the thumb and the heel of the hand.
\begin{itemize}
    \item Pain at the base of the thumb, worsened by pinching, gripping, opening jars, turning keys, or writing.
    \item Aching pain and stiffness, often worst in the morning or after a period of rest.
    \item Swelling and tenderness over the thumb base, with the joint feeling prominent.
    \item A visible bump at the base of the thumb, sometimes called the ``shoulder sign,'' with the thumb base appearing to sag or sublux dorsally.
    \item Weakness of grip and pinch, with difficulty performing fine or forceful tasks and a tendency to drop objects.
    \item Crepitus--a clicking, grating, or catching sensation--when moving the thumb.
    \item Reduced range of motion, particularly of thumb abduction and opposition.
\end{itemize}

\section*{Differential Diagnosis}
Other causes of pain and swelling around the thumb base should be considered before attributing symptoms to arthritis.
\begin{itemize}
    \item De Quervain's tenosynovitis, in which pain is located over the radial wrist tendons and reproduced by ulnar deviation of the wrist.
    \item Scaphoid or trapezium fracture, or an old, unhealed wrist injury.
    \item Trigger thumb, which causes painful clicking and locking of the interphalangeal joint.
    \item Carpal tunnel syndrome, with numbness and tingling in the thumb and fingers rather than joint pain.
    \item A ganglion cyst at the wrist or thumb base.
    \item Inflammatory arthritis, such as rheumatoid arthritis, psoriatic arthritis, or gout, which may affect the hand symmetrically with prominent swelling.
\end{itemize}

\section*{When to Seek Medical Care}
\begin{itemize}
    \item See a doctor if thumb base pain or weakness persists beyond a few weeks or interferes with daily activities.
    \item Seek assessment if swelling, redness, or a new bump develops at the thumb base.
    \item See a doctor if you are dropping objects or struggling to hold tools, utensils, or a pen.
    \item Seek advice after a fall or injury to the thumb or wrist if pain persists, as a fracture may be present.
    \item Consult a doctor if over-the-counter pain relief is required regularly for thumb pain.
\end{itemize}

\textbf{Call 000 (or local emergency services) for:}
\begin{itemize}
    \item Sudden severe pain, swelling, or deformity of the thumb or wrist after an injury.
    \item A thumb or hand that is pale, cold, numb, or blue, suggesting impaired circulation.
    \item Severe joint pain with fever, which may indicate a septic (infected) joint.
\end{itemize}

\section*{Diagnosis and Investigations}
The diagnosis is usually made on clinical grounds, with imaging used to confirm severity and exclude other pathology.
\begin{itemize}
    \item \textbf{History:} the pattern of pain with pinching and gripping, morning stiffness, and any history of injury or repetitive activity.
    \item \textbf{Examination:} tenderness and swelling over the thumb base, a prominent ``shoulder sign,'' crepitus on movement, and reduced pinch and grip strength.
    \item \textbf{Provocation tests:} longitudinal compression and axial grind testing of the thumb base reproduce the characteristic pain.
    \item \textbf{X-rays:} show joint space narrowing, subchondral sclerosis, osteophytes (bone spurs), and, in advanced cases, subluxation of the joint; findings are graded from 1 to 4.
    \item \textbf{MRI or ultrasound:} useful when the diagnosis is uncertain, to assess tendons, ligaments, or early cartilage damage.
    \item \textbf{Diagnostic injection:} injection of local anaesthetic into the CMC joint that abolishes the pain confirms the joint as the source.
\end{itemize}

\section*{Management}
Treatment follows a stepwise approach, beginning with simple conservative measures and reserving surgery for severe, refractory cases.
\begin{itemize}
    \item \textbf{Splinting:} a thumb spica splint worn during painful or heavy activities rests the joint and relieves symptoms.
    \item \textbf{Activity modification and joint protection:} using larger grips, both hands, and adapted tools to reduce load on the thumb base.
    \item \textbf{Medication:} paracetamol or non-steroidal anti-inflammatory drugs (such as ibuprofen) given topically or orally for pain and inflammation.
    \item \textbf{Hand therapy:} a hand therapist or occupational therapist provides strengthening, stretching, and ergonomic strategies.
    \item \textbf{Heat and cold:} warm packs ease morning stiffness, while cold packs reduce pain and swelling in flare-ups.
    \item \textbf{Corticosteroid injection:} an image-guided or palpation-guided injection into the joint can provide relief lasting weeks to months.
    \item \textbf{Surgery:} for severe pain, most commonly trapeziectomy (removal of the trapezium), sometimes combined with ligament reconstruction and tendon interposition.
\end{itemize}

\section*{Complications}
\begin{itemize}
    \item Progressive loss of thumb motion, grip strength, and pinch strength if left untreated.
    \item Increasing difficulty with everyday tasks such as writing, dressing, cooking, and opening packaging.
    \item A more pronounced bump and subluxation deformity of the thumb base over time.
    \item Chronic pain and reduced quality of life and independence.
    \item Side effects of long-term pain medication, including gastrointestinal and renal effects with NSAIDs.
    \item Surgical complications, including infection, nerve injury, stiffness, scar tenderness, and occasionally persistent or recurrent pain.
\end{itemize}

\section*{Prognosis and Outcomes}
CMC arthritis is a slowly progressive condition, but its course varies considerably between individuals. Many people remain comfortable for years with splinting, injections, and activity modification, and some experience spontaneous improvement. When conservative measures fail, surgery reliably relieves pain in the majority of patients, although complete restoration of normal strength is not guaranteed, and recovery of pinch power typically takes several months. Early treatment and consistent joint protection slow progression and preserve hand function, and most people ultimately retain a useful, functional thumb.

\section*{Prevention and Self-Care}
\begin{itemize}
    \item Use joint-protection principles: favour larger joints, use both hands, and use grip aids for jars and taps.
    \item Take regular breaks during repetitive pinching, gripping, or twisting tasks.
    \item Keep thumb and hand muscles strong with gentle exercises advised by a hand therapist.
    \item Maintain a healthy weight and remain physically active to support overall joint health.
    \item Wear a supportive thumb splint for heavy or repetitive activities.
    \item Use heat and cold to manage symptoms, and take over-the-counter pain relief as directed.
    \item Avoid prolonged forceful pinching, such as tight grips on heavy tools, pens, or jar lids.
    \item Seek prompt medical care after thumb or wrist injuries so that they heal properly and do not accelerate arthritis.
\end{itemize}

\section*{Sources and Further Reading}
The following sources provide reliable, current information on thumb base arthritis for both healthcare professionals and patients.
\begin{itemize}
    \item National Health Service (NHS). \textit{Thumb arthritis: symptoms and treatment.} https://www.nhs.uk/conditions/osteoarthritis/
    \item Mayo Clinic. \textit{Thumb arthritis: diagnosis and treatment.} https://www.mayoclinic.org/diseases-conditions/thumb-arthritis/
    \item MSD Manual. \textit{Osteoarthritis (OA).} https://www.msdmanuals.com/professional/musculoskeletal-and-connective-tissue-disorders/arthritis/osteoarthritis
    \item MedlinePlus. \textit{Thumb arthritis.} https://medlineplus.gov/ency/article/000423.htm
    \item Versus Arthritis. \textit{Osteoarthritis information and support.} https://www.versusarthritis.org/
    \item American Academy of Orthopaedic Surgeons (AAOS). \textit{Hand osteoarthritis.} https://www.aaos.org/
\end{itemize}
